
Temporal ordering, layer-wise mechanism, forgetting stability, and gradient-aware intervention are tested through controlled ablation experiments on CMNIST benchmark. Multi-dataset validation (Waterbirds, CelebA, NICO++) is planned future work due to manual setup requirements.

\subsubsection{Datasets}

\textbf{CMNIST (Colored MNIST).} Canonical spurious correlation benchmark with controlled color-label bias. Training set: 60,000 images, color biased with label (75\% correlation---digit 0 $\rightarrow$ red, digit 1 $\rightarrow$ green). Test set: 10,000 images with reversed bias (25\% correlation) to measure worst-group accuracy. Spurious feature: digit color (low-level visual). Core feature: digit shape (edges, stroke patterns).

\textbf{Preprocessing.} Resize to $224\times 224$ for ResNet compatibility, normalize to ImageNet statistics. Ablation variants: (1) Spurious-only---Gaussian blur $\sigma =3$ removes shape edges, preserves color. (2) Core-only---grayscale conversion removes color, preserves shape. (3) Baseline---original colored sharp digits.

\subsubsection{Models and Training}

\textbf{Architecture.} ResNet-18 (pretrained ImageNet weights), final FC layer replaced with 10-class classification head. Total parameters: 11.2M.

\textbf{Optimizer.} SGD with momentum 0.9, learning rate 0.01 (constant for proof-of-concept), weight decay $10^{-4}$, batch size 256.

\textbf{Training.} 30 epochs per variant (spurious-only, core-only, baseline) to observe post-convergence behavior. Convergence tracked via per-epoch accuracy. Convergence criterion: accuracy $\geq$ 90\% for 3 consecutive epochs.

\textbf{Statistical validation.} 10 random seeds (seeds 0-9) per experiment for paired t-tests. Proof-of-concept results use seed 0 only; full 10-seed validation launched but not completed.

\subsubsection{Experimental Questions}

\textbf{Q1:} Do spurious features converge $\geq 2$ epochs earlier than core features?  
\textbf{Method:} Train spurious-only, core-only, baseline variants on CMNIST. Measure $E_s, E_c$ (convergence epochs), compute temporal gap $\Delta = E_c - E_s$.  
\textbf{Success:} $\Delta \geq 2$ epochs.

\textbf{Q2:} Is temporal gap driven by layer-wise feature hierarchy?  
\textbf{Method:} Compute neuron-spurious correlation $\rho_j$ per neuron from ablation training activations. Aggregate by layer, test early > late layers.  
\textbf{Success:} $\bar{\rho}_{\text{early}} > \bar{\rho}_{\text{late}}, p < 0.05$ via independent t-test.

\textbf{Q3:} Do spurious features exhibit lower forgetting rate?  
\textbf{Method:} Track prediction flips per training example across 30 epochs. Compute forgetting rate $F$ (events/sample) for spurious-only vs core-only.  
\textbf{Success:} $F_{\text{spurious}} < F_{\text{core}}$.

\textbf{Q4:} Can gradient-aware training match JTT worst-group accuracy?
\textbf{Method:} Use CMNIST $\rho_j$ values to modulate per-neuron learning rates on Waterbirds. Compare worst-group accuracy to JTT baseline (86\%).
\textbf{Success:} WG-Acc $\geq 85\%$.

\subsubsection{Baselines}

\textbf{ERM (Empirical Risk Minimization).} Standard training without debiasing. Expected worst-group accuracy approximately 40\% on Waterbirds.

\textbf{JTT (Just Train Twice).} State-of-the-art reweighting method. Expected worst-group accuracy approximately 86\% on Waterbirds.

\subsubsection{Evaluation Metrics}

\textbf{Temporal gap $\Delta$.} $E_c - E_s$ in epochs. Primary metric for Q1.

\textbf{Layer-wise $\rho_j$ gradient.} Mean neuron-spurious correlation difference between early (conv1, layer1) and late (layer3, layer4) layers. Mechanism validation for Q2.

\textbf{Forgetting rate $F$.} Average forgetting events per training example. Stability metric for Q3.

\textbf{Worst-group accuracy (WG-Acc).} Accuracy on minority group. Intervention metric for Q4.
